% ===================== RESTAURAR ESTILO ANTES DEL FRONTMATTER =====================
\pagestyle{corporativo}

% ===================== FRONTMATTER CON MEMBRETE FORZADO =====================
\begin{frontmatter}

	\title{Informe de investigación del sistema de ventilación y presurización del laboratorio}

	\author[inst1]{J.Arboleda\corref{cor1}}
	\ead{proyectos2@dmlsas.com}

	\author[inst1]{F. Navia}
	\ead{fernando.navia@dmlsas.com}

	\author[inst2]{H. Rosero}
	\ead{herminsul.rosero@dmlsas.com}

	\cortext[cor1]{Autor correspondiente}
	\address[inst1]{DML Ingenieros consultores, Cali, Colombia, 760025}
	\address[inst2]{Especialidad Mecánica y Procesos, Cali, Colombia}

	% ===================== ABSTRACT CON MEMBRETE GARANTIZADO =====================
	\begin{abstract}
		\thispagestyle{corporativo}
		El presente informe documenta la investigación de mercado y la selección preliminar de los equipos y componentes del sistema de ventilación por presurización positiva del laboratorio de análisis industrial de BRINSA en Cajicá (Cundinamarca, 2 558 msnm): ventilador de impulsión directa, etapa de filtración MERV 13-14, rejillas de exfiltración con malla anti-insectos, damper de alivio barométrico e instrumentación de presión diferencial, con set-point de +25 Pa y mínimo admisible de +12.5 Pa.

		La investigación establece las condiciones ambientales de diseño del sitio (presión atmosférica 74.1 kPa, densidad del aire 0.88 kg/m$^3$) y cuantifica la corrección de densidad por altitud (factor $k$ = 0.733), que gobierna la selección del ventilador en el punto equivalente de catálogo 3 840 m$^3$/h a 260 Pa. Se presenta el análisis comparativo de candidatos comerciales por componente (Greenheck, New York Blower, Sodeca, Camfil, AAF, Ruskin, Dwyer, entre otros), el criterio de materiales para el ambiente corrosivo de la planta de hipoclorito de calcio (PRFV viniléster, acero inoxidable 316, recubrimientos epóxicos) y la selección recomendada. La filtración HEPA se descarta expresamente por no tratarse de un sistema de biocontención.
	\end{abstract}

	\begin{keyword}
		ventilación de laboratorios \sep presurización positiva \sep corrección de densidad por altitud \sep ambiente corrosivo \sep filtración MERV \sep damper barométrico
	\end{keyword}

\end{frontmatter}

% ===================== RESTAURAR ESTILO DESPUÉS DEL FRONTMATTER =====================
\pagestyle{corporativo}
\thispagestyle{corporativo}
